\documentclass[12pt]{article}
\usepackage[utf8]{inputenc}
\usepackage{graphicx}
\usepackage{amsmath}
\usepackage{amssymb}
\usepackage{mathrsfs}
\usepackage{pgfornament}
\usepackage[many]{tcolorbox}
\usepackage[margin = 0.5in]{geometry}
\usepackage{collectbox}
\usepackage{mdframed}
\usepackage{xcolor}

\newcommand\textbox[1]{%
  \parbox{.333\textwidth}{#1}%
}

\linespread{1.5}

\makeatletter
\newcommand{\mybox}{%
    \collectbox{%
        \setlength{\fboxsep}{2pt}%
        \fbox{\BOXCONTENT}%
    }%
}
\makeatother

\usepackage{listings}
\usepackage{color}

\definecolor{dkgreen}{rgb}{0,0.6,0}
\definecolor{gray}{rgb}{0.5,0.5,0.5}
\definecolor{mauve}{rgb}{0.58,0,0.82}

\lstset{frame=tb,
  language=R,
  aboveskip=3mm,
  belowskip=3mm,
  showstringspaces=false,
  columns=flexible,
  basicstyle={\small\ttfamily},
  numbers=none,
  numberstyle=\tiny\color{gray},
  keywordstyle=\color{blue},
  commentstyle=\color{dkgreen},
  stringstyle=\color{mauve},
  breaklines=true,
  breakatwhitespace=true,
  tabsize=3
}
\title{STS 532 HW V}
\begin{document}

\begin{center}
\begin{large}
\textbf{MAT 523: Functions of a Real Variable I \\
Homework I \\}
\end{large}
Nolan R. H. Gagnon \\
\end{center}

\noindent\mybox{\colorbox{green}{\textbf{ 1. }}}\hspace{3mm}\hrulefill

\noindent Show that $\mathbb{C}$ (the field of complex numbers) is not an ordered field.

\vspace{5mm}

\noindent\mybox{\colorbox{lightgray}{\textbf{Solution}}}\hspace{3mm}\hrulefill


\noindent\textbf{Lemma 1A:}
\noindent $(-1)(-1)=(-1)^2 = 1.$

\noindent\textbf{Proof:}
\begin{align*}
(-1)(-1) &= (-1)(-1) + 0 \\
&= (-1)(-1)+[1 + (-1)] \\
&= [(-1)(-1)+(-1)] + 1\\
&= [(-1)(-1)+(-1)(1)] + 1 \\
&= (-1)[-1+1]+1 \\
&= (-1)(0) + 1 \\
&= 0 + 1\\
&= 1.
\end{align*}

\hfill$\blacksquare$

\noindent\textbf{Lemma 1B:} Let $F$ be a field and let $a\in F$. Then $-a = (-1)a$.

\noindent\textbf{Proof:} 

\begin{align*}
(-1)a &= (-1)a + 0 \\
&= (-1)a + [a+(-a)] \\
&= [(-1)a+a]+(-a) \\
&= [a(-1)+a(1)] + (-a) \\
&= a[(-1) + 1] + (-a) \\
&= a[0] + (-a) \\
&= 0 + (-a) \\
&= -a.
\end{align*}

\hfill$\blacksquare$

\noindent\textbf{Lemma 1C:} $1>0$.

\noindent\textbf{Proof:} By the Non-triviality Assumption, $1\neq 0$, so suppose $1<0$. Then $0-1$ is positive, and by the first positivity axiom $(0-1)(0-1)$ is positive.  But
\begin{align}
(0-1)(0-1) &= [0+(-1)][0+(-1)] \\
&= [(-1)+0][(-1)+0] \\
&= (-1)(-1) \\
&= 1
\end{align}

\noindent by Lemma 1A. By assumption $1<0$, so $(0-1)(0-1)<0$. This contradicts the fact that $(0-1)(0-1)$ is positive. Therefore, $1>0$. 

\hfill$\blacksquare$

\vspace{5mm}
\noindent\textbf{Corollary 1A:} $-1 < 0.$

\noindent\textbf{Proof:} Since $1 > 0$ by Lemma 1C, it would be a violation of the second positivity axiom if $-1 > 0$ as well. Therefore, $-1 < 0.$

\hfill$\blacksquare$

\vspace{5mm}

\noindent\textbf{Theorem:} $\mathbb{C}$ is not an ordered field.

\noindent\textbf{Proof:} Suppose $\mathbb{C}$ is an ordered field. Then an order $<$ can be defined on $\mathbb{C}$. Consider $i,0\in\mathbb{C}$. If $0<i$, then, by the first positivity axiom, it must be the case that $0<ii = i^2 = -1$. By Corollary 1A, $-1<0$, so we have a contradiction. If, on the other hand, $i<0$, then $0<-i$, by the second positivity axiom. By the first positivity axiom, it follows that 

\begin{align*}
0 &< (-i)(-i) \\
&= [(-1)i][(-1)i] \hspace{10mm} \text{(by Lemma 1B)} \\
&= [(-1)(-1)][ii] \\
&= 1i^2 \hspace{10mm} \text{(by Lemma 1A)} \\
&= i^2 \\
&= -1.
\end{align*}

\noindent Again, this is a contradiction. Therefore, $\mathbb{C}$ is not an ordered field. 
\hfill$\blacksquare$

\pagebreak

\noindent\mybox{\colorbox{green}{\textbf{2.}}}\hspace{3mm}\hrulefill

\noindent Verify the following:
\begin{flushleft}
\textcolor{blue}{(i)} For each real number $a\neq 0$, $a^2>0$. In particular, $1>0$, since $1\neq 0$ and $1 = 1^2$. \\
\textcolor{blue}{(ii)} For each positive number $a$, its multiplicative inverse $a^{-1}$ also is positive. \\
\textcolor{blue}{(iii)} If $a>b$, then $ac>bc$ if $c>0$ and $ac<bc$ if $c<0$.
\end{flushleft}

\vspace{5mm}

\noindent\mybox{\colorbox{lightgray}{\textbf{Solution}}} \hspace{3mm}\hrulefill

\noindent\textcolor{blue}{(i)} \textbf{Proof:} Let $a\in\mathbb{R}\setminus\lbrace 0 \rbrace.$ If $a>0$, then $aa = a^2 > 0$, by the first positivity axiom. If $a<0$, then $0<-a$. Thus, 

\begin{align*}
0 &< (-a)(-a) \\
&=[(-1)a][(-1)a] \hspace{10mm} \text{(by Lemma 1B)} \\
&=[(-1)(-1)][aa] \\
&= 1a^2 \hspace{10mm} \text{(by Lemma 1A)} \\
&= a^2.
\end{align*}


\hfill$\blacksquare$

\noindent\textcolor{blue}{(ii)} \textbf{Proof:} Let $a\in\mathbb{R}^{+}$. Suppose $a^{-1}<0.$ Then $0<-a^{-1}.$ Thus, 

\begin{align*}
0 &< (-a^{-1})(a) \\
&=[(-1)a^{-1}]a \hspace{10mm} \text{(by Lemma 1B)}\\
&= (-1)[aa^{-1}] \\ 
&= (-1)1 \\
&= -1 \hspace{10mm} \text{(by Lemma 1B)}
\end{align*}

\noindent By Corollary 1A, $-1<0$, so we have a contradiction. Therefore, $a^{-1}>0$.

\hfill$\blacksquare$

\noindent\textcolor{blue}{(iii)} Let $a,b,c\in\mathbb{R}$ such that $a>b$. Since $a>b$, it follows that $a-b>0.$ Suppose $c>0$. Then, by the first positivity axiom, $c(a-b)>0$. But 

\begin{align*}
c(a-b)&=c[a+(-b)] \\
&= ca + c(-b) \\
&= ac + (-b)c \\
&= ac + [(-1)b]c \hspace{10mm} \text{(by Lemma 1B)} \\
&= ac +(-1)[bc] \\
&=ac - bc.
\end{align*}

\noindent Thus, $ac-bc>0$, or $ac>bc.$ 

\noindent Now suppose $c<0$.  Then $0<-c$, so $(-c)(a-b)>0$. But 
\begin{align*}
(-c)(a-b) &= [(-1)c](a+(-b)) \\
&= (-1)[c(a+(-b))] \\
&= (-1)[ca +c(-b)] \\
&= (-1)[ca + c[(-1)b]] \hspace{10mm} \text{(by Lemma 1B)} \\
&= (-1)[ac +(-1)[cb]] \\
&= (-1)(ac) + [(-1)(-1)](bc) \\
&= (-1)(ac) + 1(bc) \hspace{10mm} \text{(by Lemma 1A)} \\
&= bc + (-1)(ac) \\
&= bc - ac.
\end{align*}

\noindent Therefore, $bc-ac>0$ or $bc>ac$.

\hfill$\blacksquare$

\pagebreak

\noindent\mybox{\colorbox{green}{\textbf{3.}}}\hspace{3mm}\hrulefill

\noindent Let $E\subseteq \mathbb{R}$ be a nonempty subset that is bounded below. Show that $E$ has a greatest lower bound and $$\inf E = -\sup \lbrace -x | x \in E\rbrace.$$

\vspace{5mm}

\noindent\mybox{\colorbox{lightgray}{\textbf{Proof}}} \hspace{3mm}\hrulefill

\noindent Let $A = \lbrace -x | x \in E \rbrace$. Since $E$ is bounded below, $A$ must be bounded above. Thus, by the Completeness Axiom, $A$ has a least upper bound $\sup A$. Since $\sup A$ is an upper bound for $A$, $-\sup A$ must be a lower bound for $E$. We show that it is the greatest lower bound for $E$. If not, then there is a lower bound $\ell$ for $E$ with $\ell > -\sup A$. If $\ell$ is a lower bound for $E$, then $-\ell$ is an upper bound for $A$ and satisfies $-\ell < \sup A$.  But this is a contradiction, since there can be no upper bound less than the supremum. Therefore, $-\sup A$ is the greatest lower bound for $E$, i.e., $$\inf E = -\sup A.$$

\hfill$\blacksquare$

\noindent\mybox{\colorbox{green}{\textbf{5.}}}\hspace{3mm}\hrulefill

\noindent Let $E\subset \mathbb{R}$ and suppose $g\in \mathbb{R}$ is an upper bound for $E$. Prove that $g$ is the least upper bound for $E$ if and only if for all $\epsilon > 0$ $E\cap[g-\epsilon, g] \neq \emptyset$, i.e., there exists $x\in E$ such that $$g-\epsilon \leq x \leq g.$$  

\vspace{5mm}

\noindent\mybox{\colorbox{lightgray}{\textbf{Proof}}} \hspace{3mm}\hrulefill

\noindent First, suppose $g = \sup E$, but there exists an $\epsilon > 0$ such that $E\cap[g-\epsilon, g] = \emptyset$. Then all $x\in E$ satisfy either $x< g - \epsilon$ or $x>g$. Clearly, no $x\in E$ can satisfy $x>g$, since $g$ is an upper bound for $E$. Therefore, $x < g - \epsilon$ for all $x\in E$. This makes $g-\epsilon$ an upper bound for $E$; however, $g - \epsilon$ is clearly smaller than $g$, which contradicts the fact that $g$ is the supremum of $E$.

\vspace{5mm}

\noindent Next, if $E\cap[g-\epsilon, g] \neq \emptyset$ for arbitrary $\epsilon > 0$, it follows that there exists an $x\in E$ such that $g-\epsilon \leq x \leq g$ for all $\epsilon > 0$. This means that no matter how small we choose $\epsilon$ to be, there will always be an element of $E$ between $g-\epsilon$ and $g$. Clearly, this means that no number smaller than $g$ can be an upper bound for $E$. Therefore, $g$ is the supremum of $E$.

\hfill$\blacksquare$

\pagebreak

\noindent\mybox{\colorbox{green}{\textbf{7.}}}\hspace{3mm}\hrulefill

\noindent Let $A$ be a bounded subset of $\mathbb{R}$. Let $B\subset \mathbb{R}$ be the
set of upper bounds for $A$. Then $\sup A = \inf B$.

\vspace{5mm}

\noindent\mybox{\colorbox{lightgray}{\textbf{Solution}}} \hspace{3mm}\hrulefill

\noindent \textbf{Lemma 7A:} The set of upper bounds for the empty set is all of $\mathbb{R}$.

\noindent \textbf{Proof:} Let $b\in\mathbb{R}$. The statement $$a\in\emptyset \text{ implies } a\leq b$$ is true because the antecedent is false. Therefore, $b$ is an upper bound for $\emptyset$. Since $b$ was an arbitrary real number, we conclude that the set of upper bounds of the empty set is all of $\mathbb{R}$.

\hfill$\blacksquare$

\vspace{5mm}

\noindent \textbf{Theorem:} Let $A$ be a bounded subset of $\mathbb{R}$. Let $B\subset \mathbb{R}$ be the set of upper bounds for $A$. Then $\sup A = \inf B$.

\noindent \textbf{Proof:} If $A = \emptyset$, then $\sup A = -\infty$. By Lemma 7A, the set of upper bounds for $A$ is $B=\mathbb{R}$. The infimum of $\mathbb{R}$ is $\inf \mathbb{R} = -\infty$. So we have $\sup A = \inf B$. Next, suppose $A\neq \emptyset$. Since $A$ is bounded, it is bounded above. By the Completeness Axiom, $A$ must, therefore, have a supremum $\sup A = b_0$. Since $b_0$ is the supremum of $A$, it follows that $b_0<b$ for all other upper bounds $b$ of $A$. That is, $b_0<b$ for all $b\in B\setminus\lbrace b_0 \rbrace$. This further implies that $B$ is bounded below and $b_0$ is a lower bound for $B$. Since $B$ is bounded below, it must have an infimum $\inf B$. We show that $\inf B = b_0$. Let $\beta$ be a lower bound for $B$. Suppose $\beta>b_0$. Then $\beta$ is not a lower bound for $B$ since $b_0\in B$ and $b_0<\beta$. This is a contradiction. Therefore, $b_0 = \inf B = \sup A$.

\hfill$\blacksquare$
\end{document}
